\begin{ZhChapter}

\chapter{Introduction}

\text 
In recent years, Unmanned Aerial Vehicle Networks (UAV Networks) have emerged as a network architecture in which multiple UAVs communicate with one another through wireless links to support information exchange and collaborative missions \cite{hayat:2016}. UAV networks have been increasingly adopted in a wide range of applications, including military operations, medical supply delivery to remote areas, e-commerce logistics, and disaster response. Owing to their ease of deployment, high mobility, and relatively low deployment cost, UAV networks can rapidly provide flexible communication services in remote or challenging environments where conventional terrestrial network infrastructure is insufficient, damaged, or difficult to deploy \cite{gu:2023}, \cite{zeng:2016}. In addition, UAVs can serve as aerial base stations or aerial relay nodes. Benefiting from the communication advantages enabled by their higher deployment altitude and Line-of-Sight (LoS) links, UAVs can extend the coverage of terrestrial networks and improve communication accessibility \cite{zeng:2016}, \cite{mozaffari:2019}. Meanwhile, UAVs have also been increasingly integrated with the Internet of Things (IoT) and Wireless Sensor Networks (WSNs) for applications such as data collection, communication relaying, and network coverage extension \cite{gu:2023}.

As the number of UAVs and the scale of UAV networks increase, conventional distributed network architectures face growing challenges in topology management, routing decisions, resource allocation, and network state monitoring. Consequently, Software-Defined Networking (SDN) has been introduced into UAV networks, giving rise to Software-Defined UAV (SD-UAV) networks \cite{abir:2023}. The fundamental concept of SDN is the separation of the control plane from the data plane, allowing network resources and forwarding behaviors to be centrally managed by an SDN controller. With a global view of the network, the SDN controller can obtain UAV node and topology information in real time and dynamically adjust routing paths, traffic forwarding policies, and resource allocation according to the current network conditions. Furthermore, centralized network visibility enables the SDN controller to perform traffic monitoring, security policy enforcement, and anomalous behavior detection. Therefore, SDN has gradually become an important technology for improving resource management and enhancing network topology security in UAV networks.

However, although the integration of SDN improves the management flexibility and programmability of UAV networks, it also introduces new cybersecurity challenges to SD-UAV networks. Since UAVs primarily communicate through wireless networks and are characterized by high mobility, frequent topology changes, limited computational and communication resources, and exposure to open environments, they are particularly vulnerable to various attacks, including eavesdropping, spoofing, jamming, and Distributed Denial-of-Service (DDoS) attacks \cite{abir:2023}, \cite{ouldrabah:2024}. In addition, although the centralized control architecture of SDN enables the SDN controller to maintain a global view of network topology and traffic information, it also makes the security of the control plane a critical factor affecting the operation of the entire UAV network. Consequently, previous studies on SD-UAV network security have primarily focused on the three fundamental security properties of confidentiality, integrity, and availability \cite{tan:2022}. In particular, some UAVs may be responsible for critical data transmission or network services. Once these critical nodes are compromised or become unavailable, communication paths may be disrupted, service requests may fail, and portions of the network may become disconnected, thereby significantly affecting the overall UAV network \cite{tan:2022}, \cite{ouldrabah:2024}. Furthermore, before launching an actual attack, an attacker may obtain routing and topology information about the UAV network through network reconnaissance or topology probing. Based on the collected information, the attacker can identify UAVs with high topological importance and select them as attack targets, allowing targeted attacks to cause more severe damage to the network than random attacks \cite{tan:2022}, \cite{hashan:2026}.

To mitigate the impact of targeted attacks on SD-UAV networks, Tan et al. \cite{tan:2022} proposed an SDN-based topology deception mechanism that exploits the SDN controller's capability to manipulate network topology information and constructs a virtual topology that differs from the actual network topology, thereby misleading attackers in identifying key UAVs. In this approach, honeypot UAVs are deployed and configured to exhibit higher topological importance in the virtual topology observed by attackers. As a result, attackers are induced to misidentify honeypot UAVs as key UAVs and prioritize them as attack targets. Since honeypot UAVs do not provide actual network services, directing attacks toward these honeypot UAVs can prevent the actual key UAVs from becoming unavailable due to direct attacks, thereby reducing the impact on overall network connectivity. By analyzing different attack-target selection strategies and changes in network connectivity following UAV node failures, the study demonstrated the effectiveness of topology deception in protecting key UAVs \cite{tan:2022}.

However, the topology deception mechanism proposed by Tan et al. \cite{tan:2022} primarily focuses on attack target selection and the impact of compromised UAV nodes on network connectivity, while the burden imposed by high-volume attacks on underlying physical network resources has not been explicitly considered. In particular, under Distributed Denial-of-Service (DDoS) attacks, an attacker may not need to directly disrupt the functionality of a UAV node. Instead, the attacker can continuously transmit a large volume of malicious packets to consume network bandwidth and forwarding resources, resulting in network congestion, packet loss, and increased communication delay, thereby degrading the availability of legitimate services \cite{ouldrabah:2024}. Therefore, even if topology deception successfully induces an attacker to select a honeypot UAV as the attack target, it does not necessarily imply that the traffic generated during the attack will have no impact on the actual UAV network.

Furthermore, the virtual topology observed by an attacker in an SD-UAV network does not necessarily correspond to the physical network topology used for actual packet transmission. When an attacker is misled by the virtual topology and directs a large volume of DDoS traffic toward a honeypot UAV, the attack packets may still need to traverse actual UAV forwarding nodes, wireless communication links, or network gateways before reaching the decoy node. If the physical forwarding paths of attack traffic and legitimate service traffic share the same network resources, a large volume of malicious packets may increase the bandwidth utilization of shared links and compete with legitimate traffic for network resources. Existing studies on link-flooding attacks have also demonstrated that an attacker does not necessarily need to directly target a service node; instead, by concentrating a large amount of traffic on specific network links, the attacker can congest those links and disrupt other legitimate communications \cite{kang:2013}. In addition, within an SDN architecture, if a large number of attack packets repeatedly trigger flow-table misses and Packet-In messages, they may further consume processing resources in both the SDN controller and the data plane \cite{wang:2015}. Therefore, under certain physical routing and resource-sharing conditions, although topology deception can reduce the probability that a key UAV is directly selected as an attack target, the attack traffic redirected toward a honeypot UAV may still indirectly affect other functional nodes and critical communication links, thereby causing collateral damage.

The above issues indicate that there may be an unexplored relationship between the deception effectiveness of topology deception and the availability of the underlying physical network. When a honeypot UAV is designed to exhibit high topological importance in the virtual topology, it may become more attractive to attackers. However, if a large volume of attack traffic is consequently concentrated on a subset of shared physical forwarding paths, the resource burden on these paths may increase and may even affect the legitimate services originally provided by key UAVs. In other words, evaluating defense effectiveness solely based on whether an attacker selects a honeypot UAV as the attack target and on the resulting network connectivity after UAV node failures may not be sufficient to fully characterize network behavior under traffic-intensive DDoS attacks. Therefore, in addition to considering the effectiveness of virtual topology manipulation in misleading attack target selection, it is also necessary to further analyze the forwarding paths of attack traffic in the physical network and their impact on limited communication resources.

Based on the above motivation, this study aims to investigate whether topology deception in SD-UAV networks may cause DDoS attack traffic to become concentrated on specific physical network resources and, furthermore, proposes a DDoS-aware topology deception mechanism that takes network resource utilization into account. The proposed approach jointly considers the deception effectiveness of the virtual topology and the resource utilization state of the underlying physical network. During virtual topology construction and attack-traffic handling, the mechanism considers the load conditions of key UAVs, wireless communication links, and forwarding resources in order to reduce resource contention between malicious and legitimate traffic. In addition, attack-traffic steering and resource isolation strategies are employed to prevent traffic attracted by honeypot UAVs from excessively consuming critical network resources. By maintaining the deception capability of topology deception while reducing the collateral damage caused by DDoS attacks, this study aims to improve service availability and network resilience in SD-UAV networks under attack.
%補貢獻
The remainder of this paper is organized as follows. Section II reviews the related work. Section III describes the proposed framework and its workflow. Section IV presents the protection mechanism for the physical UAV network. Section V presents the experimental results and analysis. Finally, Section VI concludes this paper.
\end{ZhChapter}